% =====================================================================
% CHAPTER 2: BACKGROUND AND RELATED WORK                    budget: 6 pages
% ROLE: Give the reader exactly the concepts that chapters 3-7 use:
%       (a) how hierarchy shows in text, (b) LLMs as simulated people,
%       (c) LLMs as judges, (d) how to compare generated and real text.
%       Anything not used later does not belong here.
% Literature notes on OneDrive: identifying_hierarchy.md,
%   authoritativeness-detection-text-research.md,
%   holistic-tone-authoritativeness-research.md,
%   llm-as-judge-communication-emotion-research.md,
%   llm-email-simulator-mirroring-research.md
% =====================================================================
\chapter{Background and Related Work}
\label{ch:background}

\section{Hierarchy and power in workplace language}
% budget: 2 pages

% P1 KEY: Formal rank and the power people show in text are not the same
%         thing. This thesis uses formal rank.
% SUPPORT:
%   - Bramsen et al. 2011 (Extracting Social Power Relationships).
%   - Prabhakaran et al. ("Who's (Really) the Boss?"): situational power
%     differs from the org chart.
%   - Nguyen et al. survey; Lam et al. 2018: text power and formal rank
%     often differ.
%   - Why it matters later: limitation in Ch. 8.

% P2 KEY: Research names several language markers of power. These are the
%         candidates for the Q1 measures.
% SUPPORT:
%   - Orders / controlling acts: Prabhakaran et al. (Overt Display of Power).
%   - Pronoun "you": Kacewicz et al. 2014.
%   - Hierarchy phrases: Gilbert 2012 (built on Enron).
%   - Hedges and commitment: Prabhakaran, Ganeshkumar & Rambow 2018.
%   - Politeness and gratitude: Danescu-Niculescu-Mizil et al. 2013.

% P3 KEY: Reading formal rank from one email is hard even for trained
%         models. So only a small real signal is expected.
% SUPPORT:
%   - Bramsen 2011 and Prabhakaran 2012: hierarchical power F1 34.8 on Enron.
%   - Lam et al. 2018 (Power Networks): a neural model for power relations.
%     TODO: take its reported scores from the paper.
%   - Ziems et al. (LLMs for CSS): zero-shot LLMs can classify power, with limits.

% P4 KEY: People also judge power from the whole message, not only from
%         single words. This motivates a holistic tone judge.
% SUPPORT:
%   - Körner et al. 2024: power is perceived from short texts, with consensus.
%   - Lens model: cue validity (does a cue go with the true trait?) and
%     cue utilization (does the reader use it?).
%   - Danescu-Niculescu-Mizil 2013: politeness as a human rating of the
%     whole message.

% P5 KEY: The Enron corpus is the standard data for this work, with a known
%         hierarchy.
% SUPPORT:
%   - About 500,000 emails from Enron employees (Exposé p.5).
%   - Agarwal et al. (Comprehensive Gold Standard for the Enron hierarchy).
%   - Used by most power studies above, so results can be compared.

\section{LLMs as simulated people}
% budget: 2 pages

% P6 KEY: LLMs given a persona can act as human-like agents in experiments,
%         with known caveats.
% SUPPORT:
%   - Horton et al. (Homo silicus): replicate economic experiments.
%   - Their caveats: memorization, "prompt hacking", causal inference is hard.
%   - "Using Large Language Models to Simulate Human Behavioural Experiments"
%     (Port of Mars game). TODO: check authors for the citation.

% P7 KEY: Agent simulations add memory and social settings.
% SUPPORT:
%   - Park et al. 2023 (Generative Agents): memory stream, retrieval, reflection.
%     Our persona memory follows this (Ch. 3).
%   - SOTOPIA: social goals and a holistic evaluation.
%   - Simulating Teams (VirT-Lab): team simulation and ground-truth checks.

% P8 KEY: The closest prior work studies power and workplace email
%         directly, but not against real email on the same measures.
% SUPPORT:
%   - "Do LLM Agents Mirror Socio-Cognitive Effects in Power-Asymmetric
%     Conversations": persona pairs like manager-employee.
%   - Moral Mazes in the Era of LLMs (HR Simulator): human vs. LLM email tone.
%   - Training tools (CommCoach, GLOSS) simulate workplace partners (Exposé p.7).

% P9 KEY: Simulators have known failure modes: echoing the input, agreeing
%         with the user, and repeating stock text.
% SUPPORT:
%   - Hao et al. (echo of prompt); Sharma et al. 2024 (sycophancy).
%   - Li et al. 2016 (dull, generic replies).
%   - Enron is likely in training data, so named personas could recall real
%     people. We use no names (Ch. 3).

\section{LLMs as judges}
% budget: 1.5 pages

% P10 KEY: An LLM judge scores text on a rubric. This is now common for
%          evaluating generated dialogue.
% SUPPORT:
%   - Zheng et al. 2023: pointwise scoring, MT-Bench.
%   - Gu et al. (survey): reliability has three axes, bias is one.
%   - Written score anchors (The Biased Oracle); several named dimensions
%     reported separately (ESC-Judge, HEART).
%   - Törnberg: LLMs can outperform expert coders at annotation.

% P11 KEY: Judges have known biases. The ones that matter here are
%          self-preference, length and rank.
% SUPPORT:
%   - Zheng et al. 2023: self-enhancement and verbosity bias. Advice: judge
%     from another model family.
%   - The Biased Oracle: rank bias, use more than one judge.

% P12 KEY: A judge must be checked before use: against itself, against other
%          judges, and against people.
% SUPPORT:
%   - Stability across repeated runs (voice-agent judge paper, arXiv:2608.24314).
%   - Human-coded samples as the standard check (ESC-Judge, HEART, MT-Bench).

\section{Comparing generated and real text}
% budget: 0.5 pages

% P13 KEY: Realism can be tested without a judge: by a classifier, by
%          embeddings, and by diversity.
% SUPPORT:
%   - Paired design: human and model replies to the same input (HEART).
%   - Style embeddings that ignore topic: StyleDistance (Patel et al. 2025).
%   - Self-BLEU (Zhu et al. 2018, Texygen), named in the Exposé.

\section{Research gap}
% budget: 0.5 pages

% P14 KEY: No study combines a role-based email simulator, a real-email
%          benchmark on the same measures, and a tested judge, using only
%          free models.
% SUPPORT:
%   - P8: prior work uses invented settings, no real benchmark.
%   - P11-P12: judge checks are rarely done with free, same-family judges.
%   - Leads to Q1-Q3.
